% !TEX program = pdflatex
\documentclass[11pt]{article}
\usepackage[margin=1in]{geometry}
\usepackage{amsmath,amssymb,amsthm,mathtools}
\usepackage{bm}
\usepackage{microtype}
\usepackage{enumitem}
\usepackage{booktabs}
\usepackage{graphicx}
\usepackage{hyperref}
\usepackage[nameinlink]{cleveref}
\usepackage{listings}
\usepackage[T1]{fontenc}
\usepackage[utf8]{inputenc}
\usepackage{courier}
\lstdefinestyle{mypython}{
  language=Python,
  basicstyle=\linespread{1.0}\ttfamily\small,
  columns=fullflexible,
  breaklines=true,
  frame=single,
  showstringspaces=false,
  tabsize=2,
  captionpos=b,
  keepspaces=true
}

\newtheorem{proposition}{Proposition}
\newtheorem{definition}{Definition}
\newtheorem{remark}{Remark}

\title{Certified Graph Learning via the AU--Math Triad: Measure--Compare--Equalize}
\author{Mubeen Mohammad}
\date{\today}

\begin{document}
\maketitle

\begin{abstract}
We instantiate the AU--Math triad --- Measure, Compare, Equalize --- on finite weighted graphs. Measurement is the Dirichlet energy $E[u]=\tfrac12 u^\top L u$, comparison ranks candidates by $E$ and spectral surrogates (e.g., the spectral gap $\lambda_2$), and equalization is the Euler--Lagrange condition $Lu=0$ (or descent to that state via heat flow). We give a per--update safety certificate based on the spectral gap, an auditable policy for rollback, and a fully reproducible Python reference implementation. We further formalize advanced operational layers: multiscale certificates, adaptive control (Lyapunov), uncertainty quantification, algebraic topology and sheaf extensions, optimal transport interfaces, and differentiable certification for learning architectures.
\end{abstract}

\section{Introduction: AU--Math as Operating System}
The AU--Math triad provides a minimal lens for mathematical practice:
\begin{itemize}[leftmargin=1.2em]
  \item \textbf{Measurement} $\to$ energies, norms, measures.
  \item \textbf{Comparison} $\to$ inequalities, orderings, spectral surrogates.
  \item \textbf{Equalization} $\to$ equations, invariants, Euler--Lagrange criticality.
\end{itemize}
On graphs, this becomes concrete: the Laplacian $L$ (operator) and Dirichlet energy $E$ (measurement) induce orders via spectral quantities (comparison), and equalization corresponds to solving $Lu=0$ or performing gradient steps $u\mapsto u-\alpha L u$.

\section{Reality Interface: Graph Model and Energy}
Let $G=(V,E,w)$ be a finite weighted graph with symmetric weights $w_{ij}\ge 0$. Define the adjacency $A_{ij}=w_{ij}$, degree $D_{ii}=\sum_j w_{ij}$, and unnormalized Laplacian $L=D-A$. The Dirichlet energy is
\begin{equation}
E[u]=\tfrac12\,u^\top L u \,=\, \tfrac12\sum_{(i,j)\in E} w_{ij}\,(u_i-u_j)^2.\label{eq:energy}
\end{equation}
Let $0=\lambda_1\le\lambda_2\le\dots\le\lambda_n$ denote the eigenvalues of $L$ with orthonormal eigenvectors $\{v_k\}$.

\section{Per--Step Contraction Certificate}
Let $\bar u=\tfrac{1}{|V|}\sum_i u_i$ and $u_\perp=u-\bar u\mathbf 1$. Consider the explicit heat step $u^+=u-\alpha L u$.
\begin{proposition}[Contraction on the mean-zero subspace]
If $0<\alpha<2/\lambda_{\max}$, then
\begin{equation}
\|u_\perp^+\|_2\;\le\; (1-\alpha\lambda_2)\,\|u_\perp\|_2.\label{eq:contract}
\end{equation}
\end{proposition}
\begin{proof}[Sketch]
Diagonalize $L=V\Lambda V^\top$. The update multiplies the $k$th eigen-coordinate by $1-\alpha\lambda_k$. The constant mode ($k=1$) is unchanged; all other coordinates are contracted by at most $1-\alpha\lambda_2$.
\end{proof}
\noindent \textbf{Stable steps:} $0<\alpha<2/\lambda_{\max}$. \; \textbf{Certificate policy:} enforce~\eqref{eq:contract} each step; on violation, rollback and shrink $\alpha$.

\section{Algorithm 1: Certified Heat--Flow Smoother}
Given $G$, $L$, observations $y$, steps $T$:
\begin{enumerate}[leftmargin=1.2em]
  \item Estimate/compute $\lambda_2$ and $\lambda_{\max}$; pick $\alpha\in(0,2/\lambda_{\max})$.
  \item Initialize $u_0\leftarrow y$. For $t=0,\dots,T-1$:
  \begin{enumerate}
    \item $u_{t+1}\leftarrow u_t-\alpha L u_t$.
    \item Compute $r_t=\|u_{t+1}^\perp\|/\|u_t^\perp\|$; bound $b=1-\alpha\lambda_2$.
    \item Require $r_t\le b$; else rollback to $u_t$, set $\alpha\leftarrow\alpha/2$, retry.
  \end{enumerate}
  \item Return $u_T$ and an audit log $\{E[u_t], r_t, b\}$.
\end{enumerate}

\section{Canonical Example: Ring Graph (Reproducible)}
\textbf{Setup.} $n=200$ ring with unit weights; ground truth is a one-cycle sine; observations add Gaussian noise (RNG seed~$42$). Unnormalized Laplacian.
\\\textbf{Spectrum and step.} $\lambda_2 = 9.86879\times 10^{-4}$, $\lambda_{\max}\approx 4.0$, choose $\alpha = 0.45 < 2/\lambda_{\max} = 0.5$. Theoretical bound: $1-\alpha\lambda_2 = 0.9995559$.
\\\textbf{Energy and contraction ($T=200$).} $E[u_0]=12.3518$, $E[u_T]=0.0422$, reduction factor $0.003416$. Empirical per-step contraction: mean/median/max $0.999244/0.999551/0.999554$. Certificate satisfied each step (empirical max $\le$ bound).

\section{Spectral Policy Learner (SPL)}
Define the low-frequency cone $\mathcal F_k=\operatorname{span}\{v_1,\dots,v_k\}$ and projection $\Pi_{\mathcal F_k}(z)=V_kV_k^\top z$. A safe parameter update is
\begin{equation}
\theta^+=\Pi_{\mathcal C}\big(\theta-\eta\nabla J(\theta)\big),
\end{equation}
where $\mathcal C$ enforces (i) outputs restricted to $\mathcal F_k$ (graph setting) or (ii) closed-loop LMIs for control. Per-update certificates log $(E[u_t], r_t, 1-\alpha\lambda_2)$.

\section{Foundational Mathematical Enhancements}
\subsection{Operator Generalizations}
Weighted Laplacian $L=D-W$ with $W_{ij}=\exp(-\|x_i-x_j\|^2/\sigma^2)$; normalized Laplacian $\mathcal L=D^{-1/2}LD^{-1/2}$; magnetic Laplacian $L_\theta$ (complex phases for directed graphs); connection/sheaf Laplacians for bundle-valued signals.
\subsection{Discrete Exterior Calculus (DEC)}
Extend to $k$-forms: $d_0$ (gradient), $d_1$ (curl), and codifferential $\delta=d^*$; Hodge Laplacian $\Delta_k=d_{k-1}\delta_{k-1}+\delta_k d_k$ enables curl/divergence-aware smoothing.
\subsection{Multiscale Spectral Analysis}
Graph wavelets $\psi_{t,k}(i)=\sum_{j=1}^n g(t\lambda_j)\phi_j(k)\phi_j(i)$ for bandpass $g$ and scales $t$. Cheeger-type bounds connect $\lambda_2$ to conductance $h(G)$.

\section{Advanced Operational Layers}
\subsection{Hierarchical Certificates}
Pointwise: bound local energy changes; multiscale: require contraction of wavelet coefficients across scales; temporal: bound accelerations over time to ensure smooth evolution.
\subsection{Adaptive Control Layer}
Synthesize $P\succeq 0$ with $(I-\alpha L)^\top P(I-\alpha L)-P\preceq -Q$; choose $\alpha$ via line search under the certificate. 
\subsection{Uncertainty Quantification (UQ)}
Probabilistic certificates with estimated $\hat\lambda_2$ and confidence; Bayesian spectrum estimation from observed contractions.
\subsection{Topology and Optimal Transport}
Persistent homology to preserve essential features; sheaf Laplacians for constrained consensus; Wasserstein gradient flows and entropic JKO schemes for transport-regularized dynamics.
\subsection{Differentiable Certification}
Smooth losses that penalize certificate violations and architecture-level projection to safe sets.

\section{Complexity and Scalability}
Per-step cost is one sparse matvec $\mathcal O(m)$ for $m=|E|$. Full eigendecomposition is $\mathcal O(n^3)$ (small graphs); for large sparse graphs, use Lanczos/Arnoldi (extremal eigenpairs) or power method for $\lambda_{\max}$. Crude bounds: Gershgorin / max degree for $\lambda_{\max}$; Cheeger-type for $\lambda_2$. Frequency projection can leverage randomized SVD or Nystr"om with near-linear cost.

\section{Implementation Architecture}
A modular certification stack composes spectral, multiscale, temporal, probabilistic, and thermodynamic certifiers. A cross-domain \emph{CertifiedSystem} abstraction orchestrates dynamics, certificates, and fail-safes.

\appendix
\section{Appendix A: Reproducible Python (Ring Graph + Certificate)}
\lstset{style=mypython}
\begin{lstlisting}[language=Python,caption={Reference implementation with per-step certificate and rollback.}]
import numpy as np
import numpy.linalg as LA
import matplotlib.pyplot as plt

def ring_graph_laplacian(n: int):
    A = np.zeros((n, n), dtype=float)
    for i in range(n):
        A[i, (i - 1) % n] = 1.0
        A[i, (i + 1) % n] = 1.0
    D = np.diag(A.sum(axis=1))
    L = D - A
    return L, A, D

def dirichlet_energy(u: np.ndarray, L: np.ndarray) -> float:
    return 0.5 * float(u.T @ L @ u)

def check_contraction(u_old: np.ndarray, u_new: np.ndarray, alpha: float, lambda_2: float):
    """Return (ok, actual_ratio, theoretical_bound) with robust zero-kernel handling."""
    u_old_perp = u_old - np.mean(u_old)
    u_new_perp = u_new - np.mean(u_new)
    theoretical_bound = 1.0 - alpha * lambda_2
    eps = 1e-12
    denom = LA.norm(u_old_perp)
    if denom > eps:
        actual_ratio = LA.norm(u_new_perp) / denom
        return (actual_ratio <= theoretical_bound), float(actual_ratio), float(theoretical_bound)
    else:
        is_constant = LA.norm(u_new_perp) < eps
        actual_ratio = 0.0 if is_constant else float("inf")
        return is_constant, actual_ratio, float(theoretical_bound)

def heat_flow_certified(y: np.ndarray, L: np.ndarray, alpha: float, steps: int, lambda_2: float):
    u = y.copy()
    energies = [dirichlet_energy(u, L)]
    ratios, bounds, ok_flags = [], [], []
    for _ in range(steps):
        u_old = u.copy()
        u = u - alpha * (L @ u)
        energies.append(dirichlet_energy(u, L))
        ok, r, b = check_contraction(u_old, u, alpha, lambda_2)
        ratios.append(r); bounds.append(b); ok_flags.append(ok)
        if not ok:
            u = u_old  # rollback
            alpha *= 0.5
            u = u - alpha * (L @ u)  # retry once with smaller step
            ok, r, b = check_contraction(u_old, u, alpha, lambda_2)
            ratios[-1] = r; bounds[-1] = b; ok_flags[-1] = ok
    return u, np.array(energies), np.array(ratios), np.array(bounds), bool(np.all(ok_flags))

def project_low_frequency(z: np.ndarray, L: np.ndarray, k: int) -> np.ndarray:
    vals, vecs = LA.eigh(L); V = vecs[:, :k]
    return V @ (V.T @ z)

if __name__ == "__main__":
    n = 200
    L, A, D = ring_graph_laplacian(n)
    vals = LA.eigvalsh(L)
    lambda_1, lambda_2, lambda_max = float(vals[0]), float(vals[1]), float(vals[-1])
    x = np.arange(n)
    ground_truth = np.sin(2 * np.pi * x / n)
    rng = np.random.default_rng(42)
    observed = ground_truth + 0.3 * rng.normal(size=n)
    alpha_upper = 2.0 / lambda_max
    alpha = 0.9 * alpha_upper
    steps = 200
    u, energies, ratios, bounds, all_ok = heat_flow_certified(observed, L, alpha, steps, lambda_2)
    print("n:", n)
    print("lambda_1:", lambda_1)
    print("lambda_2:", lambda_2)
    print("lambda_max:", lambda_max)
    print("alpha_upper:", alpha_upper)
    print("alpha:", alpha)
    print("E0:", energies[0])
    print("Ef:", energies[-1])
    print("reduction_factor:", energies[-1] / energies[0])
    valid = ratios[np.isfinite(ratios)]
    if valid.size:
        print("actual_ratio_mean:", float(np.mean(valid)))
        print("actual_ratio_median:", float(np.median(valid)))
        print("actual_ratio_max:", float(np.max(valid)))
    print("theoretical_bound_per_step:", float(bounds[-1]))
    print("certificate_passed_all_steps:", all_ok)
    k = 10
    projected = project_low_frequency(observed, L, k)
    print("low-frequency L2 error:", float(LA.norm(projected - ground_truth)))
    # Optional plots (avoid styling requirements)
    # plt.figure(); plt.plot(x, observed); plt.title("Observed (noisy)")
    # plt.figure(); plt.plot(x, u); plt.title("Smoothed (heat flow)")
    # plt.figure(); plt.plot(x, ground_truth); plt.title("Ground truth (sine)")
    # plt.show()
\end{lstlisting}

\section{Appendix B: Extended Modules (Stubs)}
\begin{lstlisting}[language=Python,caption={Multiscale, adaptive control, and UQ stubs for future modules.}]
import numpy as np
import numpy.linalg as LA

def pointwise_certificate(u_old, u_new, L, alpha):
    local_energy_change = np.abs(L @ (u_new - u_old))
    pointwise_bounds = alpha * np.abs(L @ (L @ u_old))
    return np.all(local_energy_change <= pointwise_bounds)

class MultiscaleCertifier:
    def __init__(self, L, scales=(10, 50, 100)):
        self.banks = [create_wavelet_basis(L, s) for s in scales]
    def check(self, u_old, u_new):
        certs = []
        for W in self.banks:
            old_c, new_c = W.T @ u_old, W.T @ u_new
            certs.append(np.linalg.norm(new_c) <= np.linalg.norm(old_c))
        return all(certs)

def create_wavelet_basis(L, scale):
    vals, vecs = LA.eigh(L)
    g = np.exp(-scale * vals)  # placeholder: heat-kernel filter
    return vecs @ np.diag(g)

class AdaptiveStepController:
    def __init__(self, L, safety_margin=0.1, certificate=None):
        self.L = L; self.margin = safety_margin; self.certificate = certificate
        self.lam_max = power_method_radius(L)
    def compute_step(self, u, grad, lambda_2_est=1e-9):
        alpha = (2.0 / self.lam_max) * (1.0 - self.margin)
        for _ in range(10):
            cand = u - alpha * grad
            ok = True
            if self.certificate is not None:
                ok, *_ = self.certificate(u, cand, alpha, lambda_2=lambda_2_est)
            if ok: return alpha
            alpha *= 0.7
        return alpha

def power_method_radius(L, iters=50, seed=0):
    rng = np.random.default_rng(seed)
    z = rng.normal(size=L.shape[0])
    z /= np.linalg.norm(z)
    for _ in range(iters):
        z = L @ z
        z /= np.linalg.norm(z)
    return float(z @ (L @ z) / (z @ z))

class BayesianSpectrum:
    def __init__(self, prior_scale=1.0):
        self.scale = prior_scale
    def posterior(self, moments):
        mu = np.mean(moments, axis=0); var = np.var(moments, axis=0)
        return mu, var
\end{lstlisting}

\section*{Reproducibility Notes}
All numeric values in the ring-graph example are produced by Appendix~A with RNG seed $42$. The certificate is logged each step and enforces rollback on violation.

\end{document}
